\subsubsection{Invariant sublattice and quadratic unit of the aggregated operator}
\label{mg02:subreticula}

Aggregation by residue classes preserves an integral structure that
specifies the preceding spectral decomposition. The initial operator is
\(M_\Pi\), obtained by averaging the multiplicative sheet of APP over
the blocks \(C_a\times C_b\). Its integral domain is \(\mathbb Z^3\),
with coordinates ordered by \(C_1,C_2,C_0\).

\begin{proposition}[Decomposition on an index-two sublattice]
Let
\[
 v_-=(1,-1,0)^{\mathsf T},\qquad
 v_+=(1,1,0)^{\mathsf T},\qquad v_0=(0,0,1)^{\mathsf T},
\]
and let \(L=\mathbb Zv_-\oplus\mathbb Zv_+\oplus\mathbb Zv_0\).
Then
\begin{equation}
 L=\{(x,y,z)\in\mathbb Z^3:x\equiv y\pmod2\},
 \qquad [\mathbb Z^3:L]=2.
 \label{mg02:indice}
\end{equation}
The sublattice \(L\) is invariant under \(M_\Pi\). In the ordered
basis \((v_-,v_+,v_0)\), the restriction has matrix
\begin{equation}
 [M_\Pi|_L]=(-1)\oplus3H,\qquad
 H=\begin{pmatrix}3&2\\4&3\end{pmatrix},\qquad \det H=1.
 \label{mg02:restriccion}
\end{equation}
\end{proposition}

\begin{proof}
The combination \(av_-+bv_++cv_0=(a+b,-a+b,c)\) has first and
second coordinates of equal parity. Conversely, for a vector with this
parity, its coordinates in the stated basis are \(a=(x-y)/2\),
\(b=(x+y)/2\) and \(c=z\), all integers. The change-of-basis matrix
and its rational inverse are
\[
 P=\begin{pmatrix}1&1&0\\-1&1&0\\0&0&1\end{pmatrix},\qquad
 P^{-1}=\begin{pmatrix}1/2&-1/2&0\\1/2&1/2&0\\0&0&1\end{pmatrix}.
\]
For this column order, \(\det P=2\), giving the index in
\eqref{mg02:indice}. The action on the three generators is
\[
 M_\Pi v_-=-v_-,\qquad
 M_\Pi v_+=9v_++12v_0,\qquad
 M_\Pi v_0=6v_++9v_0.
\]
These equalities prove both invariance and the identity
\[
 P^{-1}M_\Pi P=
 \begin{pmatrix}-1&0&0\\0&9&6\\0&12&9\end{pmatrix}
 =(-1)\oplus3H.
\]
Finally, \(\det H=9-8=1\).
\end{proof}

\begin{corollary}[Pell iteration on the rank-two submodule]
For \(n\geq0\), there are integers \(a_n,b_n\) such that
\begin{equation}
 H^n=\begin{pmatrix}a_n&b_n\\2b_n&a_n\end{pmatrix},\qquad
 a_n+b_n\sqrt2=(3+2\sqrt2)^n,
 \qquad a_n^2-2b_n^2=1.
 \label{mg02:pell}
\end{equation}
In particular, \((a_0,b_0)=(1,0)\), \((a_1,b_1)=(3,2)\) and
\((a_2,b_2)=(17,12)\). The eigenvalues of \(H\) are the conjugate
quadratic units \(3\pm2\sqrt2\).
\end{corollary}

\begin{proof}
Multiplication of matrices of the stated form corresponds to
\[
 (a,b)(c,d)=(ac+2bd,ad+bc),
\]
the multiplication law for \(a+b\sqrt2\). Hence the powers of
\(H\) satisfy \(a_{n+1}=3a_n+4b_n\), \(b_{n+1}=2a_n+3b_n\),
with the preceding initial data. The multiplicative norm of
\(3+2\sqrt2\) is one; thus \(a_n^2-2b_n^2=1\) at every degree.
The characteristic polynomial of \(H\), \(\lambda^2-6\lambda+1\),
gives the two conjugate units. The identity
\(H^{\mathsf T}\operatorname{diag}(2,-1)H=
\operatorname{diag}(2,-1)\) expresses the same conservation as an
integral bilinear form.
\end{proof}

The integral splitting therefore resides in \(L\); the same matrix
\(P\) gives a complete decomposition over \(\mathbb R^3\). The
factor \(3\) in the restriction preserves the scale of the APP
average: its rank-two iteration is \(3^nH^n\). The central block
\(B_c=\left(\begin{smallmatrix}7&2\\2&7\end{smallmatrix}\right)\),
in contrast, resides in two adjacent positions of the table and satisfies
\(B_c^{\mathsf T}\operatorname{diag}(1,-1)B_c=
45\operatorname{diag}(1,-1)\). Compression into three classes and
central restriction to two positions thus retain their respective
domains, bases and normalizations. The Pell structure proceeds from
the first operation through \eqref{mg02:restriccion}.
